\documentclass{tutodoc}

\usepackage{../preamble.cfg}

\usepackage{../main/main}


\begin{document}

%\section{La syntaxe du langage proposé}

\subsection{Les contextes pour les signes et/ou les variations}

Voici les contextes disponibles pour les signes et/ou les variations.

\medskip

\begin{itemize*}[before = \leavevmode\kern15pt, itemjoin = \kern15pt]
	\item \tdoclatexin{xvals}

	\item \tdoclatexin{signs}

	\item \tdoclatexin{vars}
\end{itemize*}


\begin{tdocnote}
	On peut utiliser au choix \tdoclatexin{signs} ou \tdoclatexin{vars} de façon isolée.
\end{tdocnote}


% --------------- %


\subsubsection{Le contexte \tdoclatexin{xvals} pour les valeurs pivots}

Se reporter à la section \ref{tns-math-functab-dsl-l3-ctxt-xvals} donnant la syntaxe attendue ; ce qui échange ici est que les valeurs données sont des valeurs pivots liées aux intervalles sur lesquels on donne des informations.

\end{document}
